% Replacement fragment for the HRES2--H2 part of
% paper/statistical_analysis_EN.tex.
% This file is separate and does not modify the original manuscript.

\subsection{HRES2--H2 Engineering Case}
\label{sec:statistical-hres2-en-revision}

The HRES2--H2 experiment evaluates the sizing and rule-based dispatch of a
grid-connected wind--photovoltaic--electrolyzer--battery (WPEB) system at
Baotou, Inner Mongolia, China ($41.70^\circ\mathrm{N}$,
$110.43^\circ\mathrm{E}$). The current implementation calculates annual
green-hydrogen production from the electrolyzer load, but it does not model a
minimum hydrogen-demand target, a hydrogen storage tank, or a fuel-cell
conversion subsystem. The case study is therefore interpreted as WPEB sizing
with annual hydrogen production, rather than as a complete hydrogen-storage
microgrid.

The comparison includes ten standalone metaheuristics: PSO, GWO, WOA, EHO,
ACO, ABC, ILS, SA, TS, and VNS. Each method was evaluated in 31 stochastic
optimization runs with a maximum budget of 1,000 iterations. Each objective
evaluation simulates 8,760 hourly dispatch steps over one fixed synthetic
weather profile. Thus, the experiment measures optimization stochasticity on
a common scenario; it is not a multi-year meteorological Monte Carlo study.

The inferential protocol comprises Shapiro--Wilk normality testing, the Friedman omnibus test, and pairwise Wilcoxon signed-rank tests. The pairwise values reported here are the unadjusted Wilcoxon $p$-values generated by the current implementation; no Holm correction is applied. The 31 corresponding executions were treated as paired runs for the inferential comparison.
The archived Mann--Whitney results are also available as an independent-sample
sensitivity check.

The Shapiro--Wilk test rejected normality for both proposed LCOE distributions
($p<10^{-11}$). Friedman tests detected significant differences among the
eleven methods in both comparisons: $\chi_F^2=129.1662$ for DTW
($p=6.910283\times10^{-23}$) and $\chi_F^2=123.4738$ for DDTW
($p=9.966826\times10^{-22}$). The proposed variants obtained mean ranks of
1.71 and 1.95, respectively.

All ten standalone algorithms had unadjusted Wilcoxon $p<0.001$ against the
corresponding proposed control and a higher mean rank. Consequently, the
archived values support a $10/0/0$ win/tie/loss record at the unadjusted
$\alpha=0.05$ level. The independent Mann--Whitney sensitivity test also
returned $p<0.05$ for all ten baseline comparisons in both variants. These
results support the observed LCOE advantage on this fixed engineering
scenario, without claiming multiplicity-adjusted inference.

The direct row-wise DTW--DDTW comparison yielded no significant LCOE
difference (Wilcoxon $p=0.593$). The difference in archived means was only
$7.42\times10^{-7}$ CNY/kWh, below the six-decimal reporting precision. The
available evidence therefore supports comparable economic performance of the
two alignment modalities in this case; it does not establish superiority of
either DTW or DDTW.

\begin{table}[p]
\centering
\scriptsize
\caption{Archived HRES2--H2 LCOE inferential comparison over 31 runs. The
pairwise values are unadjusted Wilcoxon signed-rank $p$-values.}
\label{tab:hres2-pvalues-revision-en}
\resizebox{\textwidth}{!}{\begin{tabular}{lrrrrrr}
\toprule
& \multicolumn{3}{c}{DTW} & \multicolumn{3}{c}{DDTW} \\
Algorithm & Mean LCOE & Mean rank & Raw $p$ & Mean LCOE & Mean rank & Raw $p$ \\
\midrule
Hybrid control & 0.267160 & 1.71 & -- & 0.267160 & 1.95 & -- \\
GWO & 0.268066 & 4.00 & $1.3812\times10^{-6}$ & 0.268395 & 4.06 & $1.6196\times10^{-6}$ \\
VNS & 0.267806 & 5.19 & $9.3132\times10^{-10}$ & 0.267806 & 5.00 & $1.3812\times10^{-6}$ \\
TS  & 0.267188 & 5.35 & $4.6566\times10^{-9}$ & 0.267184 & 5.52 & $9.3132\times10^{-10}$ \\
ACO & 0.273288 & 5.58 & $1.7539\times10^{-5}$ & 0.273617 & 5.68 & $4.1782\times10^{-4}$ \\
ILS & 0.267184 & 5.58 & $9.3132\times10^{-10}$ & 0.267183 & 5.58 & $3.4552\times10^{-7}$ \\
PSO & 0.277289 & 6.00 & $3.6279\times10^{-4}$ & 0.280068 & 7.00 & $8.0432\times10^{-5}$ \\
WOA & 0.277227 & 7.26 & $1.6455\times10^{-5}$ & 0.276477 & 5.97 & $2.9334\times10^{-4}$ \\
EHO & 0.277728 & 7.84 & $6.1660\times10^{-6}$ & 0.278524 & 8.06 & $2.2878\times10^{-6}$ \\
ABC & 0.272971 & 7.84 & $1.1714\times10^{-6}$ & 0.270993 & 7.34 & $1.1714\times10^{-6}$ \\
SA  & 0.286431 & 9.65 & $9.3132\times10^{-10}$ & 0.290062 & 9.84 & $9.3132\times10^{-10}$ \\
\bottomrule
\end{tabular}}
\end{table}

The inferential comparison intentionally reports only LCOE for the standalone
methods. The archived standalone output does not contain per-run LCOH, AGSR,
feasibility, hydrogen-production, or component-sizing records. Those metrics
should not be presented as independently auditable baseline comparisons until
the corresponding per-run files are archived.

For a stronger engineering validation, the experiment should be repeated over
several hourly weather-year profiles. Each profile should be passed to
\texttt{HRES2Function(df\_year=...)}, with all methods evaluated under the same
scenario and a documented randomization protocol. The resulting per-year and
pooled distributions should include LCOE, LCOH, AGSR, hydrogen production,
feasibility, and component sizing. Until then, the present evidence should be
described as a controlled synthetic-scenario study.
